\begin{figure}
\centering
\begin{tikzpicture}[x=7mm,y=5mm]
\draw[cfg edge] (-0.5,0) -- (7,0) node[right] {$i$};
\draw[cfg edge] (0,-0.5) -- (0,10) node[above] {$j$};
\draw[dashed] (0,3) rectangle (5,8);
\draw[semithick] (0,3) -- (5,8);
\foreach \x in {0,1,2,3,4,5} {
  \fill (\x,{\x+3}) circle (2pt);
}
\node[font=\small,anchor=west] at (5.6,6.6) {$j-i=3$};
\node[font=\small,anchor=west] at (5.6,3.5) {interval box};
\node[below] at (5,0) {$5$};
\node[left] at (0,3) {$3$};
\node[left] at (0,8) {$8$};
\draw (5,3) circle (3pt);
\node[font=\small,anchor=west] at (5.6,1.5) {$(5,3)$ is spurious};
\end{tikzpicture}
\caption{The six filled points satisfy the fixed difference. Independent intervals also admit points off the diagonal, including the open point $(5,3)$.}
\label{fig:ch43-relation}
\end{figure}
